\chapter{Összefoglalás}

A szakdolgozatban egy olyan webalkalmazás megtervezése és megírása volt a célom, amely lehetővé teszi XSD sémák alapján dinamikus űrlapok létrehozását, és a kitöltött adatokból XML dokumentumok előállítását. A fejlesztés során célom volt, hogy a rendszer a felhasználók számára egyszerűen használható felületet biztosítson, miközben az űrlapok felépítése és az adatbevitel szabályai az eredeti sémához igazodjanak.

A kitűzött célt úgy gondolom, hogy sikerült megvalósítanom. Az elkészült alkalmazásban lehetőség van a felhasználónak regisztrációra és bejelentkezésre, XSD sémák feltöltésére és szerkesztésére, és az azokból előállított dinamikus űrlapok létrehozására. A program külön felületen ad lehetőséget a sémák készítésére és az űrlapok kitöltésére, ezért azok egyszerűen elérhetőek és használhatóak a felhasználók számára. A kitöltött adatokból az alkalmazás XML dokumentumot generál, amely később letölthető és felhasználható.

A fejlesztés legösszetettebb része az XSD sémák feldolgozása és az azok alapján történő űrlapgenerálás megvalósítása volt. Nagy kihívás volt, hogy az XSD nem csupán egyszerű adattípusokat, hanem egymásba ágyazott elemeket, attribútumokat, ismétlődő elemeket és különböző értékkorlátozásokat is tartalmazhat. Ezeket úgy kellett feldolgozni, hogy az alkalmazás képes legyen a megfelelő űrlapmezőket létrehozni, és a lehető legtöbb megkötést figyelembe vegye. A fejlesztés során egy olyan feldolgozási megoldást hoztam létre, amely többféle egyszerű és összetett adatszerkezet kezelését is lehetővé teszi.

Mindezek ellenére a fejlesztés során világossá vált, hogy a rendelkezésre álló idő alatt nem lesz lehetőségem ahhoz, hogy az XSD teljes funkcionalitását támogassam. Több szerkezeti elem és típusmeghatározás kezelése korlátozott maradt. Emiatt a program jövőbeli továbbfejlesztésének egyik legfontosabb feladata ezeknek a hiányosságoknak a kiküszöbölése. Az XSD-hez tartozó elmaradásokon kívül, egy másik lehetőséget jelentene az űrlapok megjelenésének egyénileg való testreszabása.